\subsection{弱空间中的完备凸锥}

引理 1. --- 设 E 是一个 Hausdorff 弱空间，C 是 E 中以 0 为顶点的真锥，且它关于由 E 的一致结构所诱导的一致结构是完备的。则 E 上每个连续线性型都是两个在 E 上连续且在 C 上为正的线性型之差。

设 \(E'\) 为 E 的对偶，F 为 \(E'\) 的代数对偶，赋拓扑 \(\sigma(F, E')\)。设 \(H = C^{\circ} - C^{\circ}\) 为 \(E'\) 中由在 E 上连续且在 C 上为正的线性型之差所构成的向量子空间（II, p.~44, 命题 4）。只需证明 H 在 F 中的正交是 \(\{0\}\)（II, p.~41, 例 1）。于是设 \(a \in F\) 与 H 正交；由于 a 与 \(C^{\circ}\) 正交，它必属于 C 在 F 中的双极集。但 E 可等同于 F 的一个子空间，且由于 C 完备，从而在 F 中闭，故有 \(a \in C\)（II, p.~44, 定理 1）。类似地，a 与 \(-C^{\circ}\) 正交，因此 \(a \in -C\)。由于 C 是真锥，得 a = 0。

命题 11. --- 设 E 是一个 Hausdorff 弱空间，C 是 E 中以 0 为顶点的真凸锥，且它关于由 E 的一致结构所诱导的一致结构是完备的。

则存在一个集合 I 与一个 E 到乘积空间 \(R^{I}\) 中的连续线性映射 u，它具有下列性质：

\begin{enumerate}
\def\labelenumi{\alph{enumi})}
\item
  \(u\) 是 \(\mathbf{C}\) 到 \(u(\mathbf{C})\) 上的同构，两者分别赋以由 \(\mathbf{E}\) 的一致结构与 \(\mathbf{R}^{\mathrm{I}}\) 的一致结构诱导的一致结构。
\item
  我们有 \(u(\mathbf{C}) \subset \mathbf{R}_{+}^{\mathrm{I}}\)。
\end{enumerate}

此外，若 E 的一致结构在 C 上诱导的一致结构是可度量化的，则可取 I = N。

设 \((f_{\iota})_{\iota \in I}\) 为 E 上一族连续线性型，使得形如 \((x, y) \mapsto |f_{\iota}(x - y)|\) 的伪度量在 \(\mathbf{C} \times \mathbf{C}\) 上的有限和定义 C 的一致结构。（若该结构可度量化，则可取 \(\mathbf{I} = \mathbf{N}\)。）由引理 1，我们还可假设每个 \(f_{\iota}\) 在 C 上为正。设 \(u\) 为线性映射 \(x \mapsto (f_{\iota}(x))_{\iota \in I}\)，它把 E 映到 \(\mathbf{R}^{\mathrm{I}}\) 中。显然 \(u\) 连续且 \(u(\mathbf{C}) \subset \mathbf{R}_{+}^{\mathrm{I}}\)。限制 \(u|\mathbf{C}\) 是 C 到 \(u(\mathbf{C})\) 上的满的一致连续映射。此外，若 C 中的 \(x, y\) 使 \(f_{\iota}(x) = f_{\iota}(y)\) 对所有 \(\iota \in I\) 成立，则由于 C 的一致结构是 Hausdorff 的，有 \(x = y\)；因此 \(u|\mathbf{C}\) 是双射。最后，若 W 是 C 的一致结构的一个一致邻域，则存在 I 的有限子集 J 与数 \(\varepsilon > 0\)，使得关系式 \(|f_{\iota}(x) - f_{\iota}(y)| \leqslant \varepsilon\)（对 \(\iota \in J\)）蕴含 \((x, y) \in W\)；因此 \(u|\mathbf{C}\) 是 C 到 \(u(\mathbf{C})\) 上的、关于所考虑的一致结构的同构。

推论 1. --- 设 E 是一个 Hausdorff 弱空间，C 是 E 中以 0 为顶点的真凸锥，且它关于由 E 的一致结构所诱导的一致结构是完备的。则映射 \((x, y) \mapsto x + y\)（\(C \times C\) 到 C 中）是逆紧的。

由于命题 11，我们可以假设 \(\mathbf{E} = \mathbf{R}^{\mathbf{I}}\) 且 \(\mathbf{C} = \mathbf{R}_{+}^{\mathbf{I}}\)（GT, I, § 10.1, 推论 1 与 4）。但这时映射 \((x, y) \mapsto x + y\)（\(\mathbf{C} \times \mathbf{C}\) 到 \(\mathbf{C}\) 中）写作 \(((\xi_{\mathfrak{i}}), (\eta_{\mathfrak{i}})) \mapsto (\xi_{\mathfrak{i}} + \eta_{\mathfrak{i}})\)，于是我们只需证明连续映射 \(f: (\xi, \eta) \mapsto \xi + \eta\)（\(\mathbf{R}_{+} \times \mathbf{R}_{+}\) 到 \(\mathbf{R}_{+}\) 中）是逆紧的（GT, I, § 10.1, 推论 3）。现在，对所有 \(\zeta \in \mathbf{R}_{+}\)，我们看到 \(f(\zeta)\) 是诸对 \((\xi, \zeta - \xi)\)（满足 \(0 \leqslant \xi \leqslant \zeta\)）所成之集，因此 \(f\) 对区间 \([0, \zeta]\) 的逆像是 \((\xi, \eta) \in \mathbf{R}_{+} \times \mathbf{R}_{+}\) 中满足 \(\xi + \eta \leqslant \zeta\) 者所成之集，它是紧的。应用（GT, I, § 10.3, 命题 7）即得结论。

推论 2. --- 设 E 是一个 Hausdorff 弱空间，C 是 E 中以 0 为顶点的真凸锥，且它关于由 E 的一致结构所诱导的一致结构是完备的。

\begin{enumerate}
\def\labelenumi{(\roman{enumi})}
\item
  对每一点 \(a\)（属于 \(\mathbf{E}\)），交 \(\mathbf{C} \cap (a - \mathbf{C})\) 是紧的。
\item
  设 \(\mathbf{A},\mathbf{B}\) 是 C 中两个闭集。则 \(\mathbf{A} + \mathbf{B}\) 是 C 中的闭集。
\item
  由推论 1 及 GT, I, § 10.2, 定理 1, b)，\((x, y) \in \mathbf{C} \times \mathbf{C}\) 中满足 \(x + y = a\) 的那些所成之集是紧的。而这个集也就是诸 \((x, a - x)\)（其中 \(x \in \mathbf{C} \cap (a - \mathbf{C})\)）所成之集，这就证明了 (i)。
\item
  若 A 与 B 在 C 中闭，则 \(A \times B\) 在 \(C \times C\) 中闭，于是由推论 1 与 GT, I, § 10.1, 命题 1，\(A + B\) 在 C 中闭。
\end{enumerate}

